\section{SCENARIO \#7: ''New Market'': May 13-18, 1864
}


\begin{scenariodata}
\textbf{GAME LENGTH:} 6 Turns\\
\end{scenariodata}

\begin{scenariodata}
(WEATHER: Turns 1-3; May 1-15 Column; Turns 4-6: May 16-31\\
\end{scenariodata}

Column.)

\begin{scenariodata}
\textbf{FIRST PHASE:} Union\\
\end{scenariodata}

\begin{scenariodata}
\textbf{UNION:}\\
\end{scenariodata}

\begin{scenariodata}
\textbf{RESOURCE FACTOR:} 4\\
\textbf{INITIAL PLACEMENT:}\\
\end{scenariodata}

\begin{scenariodata}
B\&O : [101, 2A, 1C], (4F, 18)\\
Hex M-29 (Woodstock) : 101, 4A, 6C, 1H,\\
\end{scenariodata}

\begin{scenariodata}
\textbf{and adjacent :} (Sigel, 1F, 2S, 2W)\\
\end{scenariodata}

\begin{scenariodata}
Hex N-20 (Winchester) : 1, (1F, 1S, 1W)\\
\end{scenariodata}

\begin{scenariodata}
Hex S-28 (Front Royal) : 21, (1F, 1S, 17W)\\
and adjacent\\
\end{scenariodata}

\begin{scenariodata}
Hex R-31 21C\\
\end{scenariodata}

\begin{scenariodata}
Hex M-26 (Strasburg) : 21, (1F,1S,1W)\\
and adjacent\\
\end{scenariodata}

\begin{scenariodata}
\textbf{WITHDRAWAL HEXES:}\\
N-7, \$-10, Z-13, CC-15 through CC-33\\
(SUPPLY/WAGON ENTRY HEXES:\\
\end{scenariodata}

\begin{scenariodata}
N-7, S-10, Z-13, A-11, CC-15 through CC-33)\\
\end{scenariodata}

\begin{scenariodata}
No Additional Fortification counters.\\
\end{scenariodata}

\begin{scenariodata}
\textbf{CONFEDERATES:}\\
\end{scenariodata}

\begin{scenariodata}
\textbf{RESOURCE FACTOR:} 3\\
\textbf{INITIAL PLACEMENT:}\\
Hex O-51 (Staunton) : 91, 3A,\\
\end{scenariodata}

\begin{scenariodata}
\textbf{and adjacent :} (Breckinridge, 1F, 25, 2W)\\
Hex 0-35\\
\end{scenariodata}

\begin{scenariodata}
and adjacent > 21, (1F, 1S, 1W)\\
\end{scenariodata}

\begin{scenariodata}
Within 3 hexes of\\
\end{scenariodata}

\begin{scenariodata}
Hex 0-35 :2C,1H\\
\end{scenariodata}

\begin{scenariodata}
Hex H-27\\
\end{scenariodata}

\begin{scenariodata}
\textbf{and adjacent :} 2c\\
\end{scenariodata}

\begin{scenariodata}
(Partisans: McNeill, Gilmor, Mosby)\\
\end{scenariodata}

\begin{scenariodata}
\textbf{WITHDRAWAL HEXES:}\\
\end{scenariodata}

\begin{scenariodata}
CC-40 through CC-58\\
(SUPPLY/WAGON ENTRY HEXES:\\
CC-40 through CC-58, 0-51, M-50)\\
\end{scenariodata}

Three additional Fortification counters available.

\begin{scenariodata}
(SPECIAL: Place a Destruction counter on bridge M-32/N-33.)\\
\end{scenariodata}

\begin{scenariodata}
\textbf{VICTORY CONDITIONS:} AI! other results are a draw.\\
\end{scenariodata}

\begin{scenariodata}
\textbf{ADVANCED GAME:} Union wins if they have 10 or more lead in total\\
\end{scenariodata}

\subsection*{Victory Points,}

\begin{scenariodata}
Confederates win if they have 15 or more lead in\\
\end{scenariodata}

total Victory Points.

\begin{scenariodata}
\textbf{BASIC GAME:} Same as the Advanced Game, only "'spot'' the Con-\\
\end{scenariodata}

federates 10 Victory Points at the start of the game.

\begin{scenariodata}
\textbf{HISTORICAL RESULTS:} Confederate victory.\\
\end{scenariodata}

\begin{scenariodata}
\textbf{PLAYTESTER'S COMMENTS:} This is one of the best-balanced sce-\\
\end{scenariodata}

narios, and is almost always settled on the last Turn of the game. The

\begin{scenariodata}
Union closeness to the southern end of the Valley gives opportunities for\\
\end{scenariodata}

cavalry raids on the Virginia Central Railroad. The Union Player should gather together as many units as possible, strike early, and force the

\begin{scenariodata}
Confederates to fight to hold onto Staunton; with a decent attrition rate,\\
\end{scenariodata}

this strategy can win. The weakness of the B\&O garrison should allow the

\begin{scenariodata}
Confederates to keep it cut with just one Partisan. The other two Partisan\\
\end{scenariodata}

units should be used in the rear of the Union army, threatening the supply wagons. This will force the Union Player to use more of his forces to guard the wagons. The Confederates can win if they can defeat small isolated union forces, or meet the main Union army in a weakened condition.

a

Sigel failed to bring his entire army into action, and the Federals were decisively beaten and driven off the field, a feat that is attributed by popular legend to a charge by the teen-aged VMI Cadet Corps. The VMI Cadets were certainly prominent in the victory, but it required much hard fighting by all Confederate units engaged to pull it off. Much of the credit for the victory should also go to the partisan leader John S. Mosby, the '"'Gray Ghost."' Mosby had begun his partisan operations in 1863, and his constant raids on Sige!'s wagon trains had been an important factor in keeping many Union units away from the critical battlefield.

Following the victory at New Market, Breckinridge led most of his command to join Lee and the Army of Northern Virginia, taking part in the Battle of Cold Harbor. While absent, however, the Union forces, in overwhelming numbers, advanced a!most unopposed up the Valley.

This fresh Union advance was commanded by David Hunter (1802-1886), a regular army officer, and his habit of writing nasty letters of criticism about his superiors did little to endear him to his fellow officers. He was particularly hated in the south for his views on slavery, and his destructive bent. Hunter took Staunton, then advanced far south to Lynchburg, his ultimate objective. There he was repulsed by Breckinridge's division and Early's If Corps of the Army of Northern Virginia, and retreated into West Virginia.

Hunter's withdrawal left the Valley open, and the Confederates advanced into the void this created, a massive raid which took the

\begin{scenariodata}
Confederates to the outskirts of Washington, D.C., and forced the\\
\end{scenariodata}

diverting of significant Union forces to save the nation's capitol. Faced by the build-up of Union forces at Washington, and the re-entry of Hunter's army into the Shenandoah Valley to his rear, the Confederates had to retreat, encumbered by their wagons, and with superior Union forces closing in on them...

